\documentclass{article}%
\usepackage[T1]{fontenc}%
\usepackage[utf8]{inputenc}%
\usepackage{lmodern}%
\usepackage{textcomp}%
\usepackage{lastpage}%
\usepackage{geometry}%
\geometry{margin=2.2cm}%
%
\usepackage{amsmath}%
\usepackage{amssymb}%
\usepackage{booktabs}%
\usepackage{longtable}%
\usepackage{array}%
\usepackage{hyperref}%
\hypersetup{hidelinks}%
\title{IME 2026 {-} Física (claude\_sonnet\_5\_5) x Poliedro\\[4pt]\large Comparação com o gabarito Poliedro Resolve\thanks{Gerado com gpt-resolve: \url{https://github.com/lgabs/gpt-resolve}}}%
\date{10/11/2026}%
%
\begin{document}%
\normalsize%
\maketitle%
\section*{Resumo}\textbf{Nota média: 10,0/10} (26 de 26 itens com a resolução correta; cada item vale o mesmo dentro da questão).\par\medskip%
\begin{longtable}{@{}c c p{0.78\textwidth}@{}}\toprule Questão & Nota & Resumo \\ \midrule
1 & 10,0 & A resolução e o gabarito do Poliedro concordam nos três itens. a) -2 m < x < 4 m. b) E\_sup = 32/3 J $\approx$ 10,67 J. c) O esboço de v$\times$t tem a mesma forma nos dois: v começa em -v0, anula{-}se no retorno (x < -2), vale +v0 em T2 = 2T1, cai até um mínimo positivo em x = 4 e depois cresce sem limite. Refiz as contas: a identidade Ep - 32/3 = -(x-4)$^2$(x+2)/3 confere, e os instantes e velocidades do esboço da resolução batem com uma integração RK4 independente. A resolução é mais detalhada que o gabarito: concavidades, inflexão em x = 2, valor de v\_min e assíntota em tempo finito. O único deslize encontrado é um erro de digitação na conta do item b do gabarito (fator 4 a mais), sem efeito no resultado final. \\
2 & 10,0 & A resolução e o gabarito do Poliedro coincidem nos cinco itens: a) 16/9; b) 3/4; c) 1/8; d) 2; e) apoastro. Refiz as contas com frações exatas e todas conferem. Também calculei a = 8r\_B/7, e = 5/9 e periastro = 32r\_B/63, e a energia bate no periastro. Na justificativa do item e, a resolução é mais completa que o gabarito, porque explica por que o ponto do choque é um apside. Há só uma pequena diferença de transcrição ('e que' em vez de 'e quando'), que não muda nenhuma resposta. \\
3 & 10,0 & A resolução está correta nos dois itens e coincide com o gabarito Poliedro: a) d\_min = 3?/5; b) $\rho$ = ?/14. Refiz as contas de forma independente e confirmei numericamente o ponto de afrouxamento (30$^\circ$ acima da horizontal), v$^2$ = 2g?/7 e v\_x$^2$ = g?/14. Também confirmei que, no trecho de lançamento oblíquo, a esfera não volta a esticar o fio nem toca a barra antes do ápice. As diferenças entre os dois textos são só de convenção do ângulo. A resolução ainda traz conferências extras (posição do ápice e consistência de energia), que estão todas corretas. \\
4 & 10,0 & A resolução e o gabarito do Poliedro chegam ao mesmo resultado: E = 7560 J. Os dois usam R\_eq = 2 $\Omega$ e potências de 50, 128 e 200 W, cada uma durante 20 s. Pela figura original, a ordem das chaves faz com que 1, 2 e 3 ramos fiquem em paralelo, um número por intervalo. Refiz as contas em Python e estão corretas. A resolução ainda acrescenta verificações certas: Millman, balanço de potência e convergência da série na escada. Não há divergência. \\
5 & 10,0 & A resolução e o gabarito do Poliedro coincidem nos dois itens. No (a), C = 22/3 N, D = 200/3 N, E = 26 N e G = 26 N. No (b), m\_máx = 8,4 kg, no limite em que C = 0. Refiz as contas com frações exatas, inclusive o equilíbrio de forças e de momentos, e os dois estão certos. O modelo é o mesmo: a viga é dividida em G, GE se apoia no dente e CG tomba em torno de D. Não encontrei erro; as dúvidas registradas são só de interpretação do enunciado. \\
6 & 10,0 & A resolução está certa e concorda com o gabarito Poliedro nos três itens: (a) 1,25 A, (b) 56,25 J e (c) 2,8125 N para a direita. Os dois usam o mesmo modelo, $\varepsilon$ = 10 V e r = 2 $\Omega$, com a fem de movimento somando{-}se a $\varepsilon$ quando v > 0, e esse sinal confere com a Figura 1 do original. Refiz as contas e não achei erro. A resolução ainda traz um balanço de potência e uma análise do sinal da fem que o gabarito não tem. \\
7 & 10,0 & A resolução e o gabarito do Poliedro coincidem nos três itens, e os dois estão certos. No item a, V(a) = kQ/(2a). No item b, V(a) = 5kQ/(6a) e V(2a) = V(3a) = kQ/(3a), com o gráfico correto. No item c, as superfícies internas ficam com -Q e as externas com +4Q/5 (casca menor) e +6Q/5 (casca maior). Conferi o enunciado no PDF original e refiz as contas com frações exatas. Não há divergência que mude a nota; só ficaram observações menores sobre hipóteses implícitas e o formato do gráfico. \\
8 & 10,0 & A resolução está correta e coincide com o gabarito Poliedro nos dois itens. Confirmei a transcrição contra a página 8 da prova. Pelas minhas contas, Q\_f = 120000 kcal = 4,8$\cdot$10\^{}8 J e a potência frigorífica é 83,3 kW. O COP de Carnot é 4, o real é 3, e a potência do compressor dá $\approx$ 27,8 kW $\approx$ 37 HP. No item b, os dois concluem que 30 HP não bastam e que a estimativa do engenheiro não procede. A resolução ainda traz verificações extras coerentes: com 30 HP só se retirariam 81\% do calor e o congelamento levaria \textasciitilde{}118,5 min. Não há erro de conta nem de física. Os únicos pontos são hipóteses implícitas comuns e redação. \\
9 & 10,0 & Questão 9: a resolução bate com o gabarito do Poliedro Resolve nos três itens, e os dois estão certos. Conferi o enunciado no PDF original (página 9) e a transcrição está fiel. Os resultados são:\newline%
a) I = I\_1{[}4 + 2$\surd$3 cos(2$\pi$a sen$\theta$/$\lambda$){]};\newline%
b) I\_máx = (4+2$\surd$3)I\_1, que é (1+$\surd$3)$^2$I\_1 $\approx$ 7,46 I\_1, em sen$\theta$ = $\lambda$/a;\newline%
c) I\_mín = (4-2$\surd$3)I\_1, que é ($\surd$3-1)$^2$I\_1 $\approx$ 0,54 I\_1, em sen$\theta$ = $\lambda$/(2a).\newline%
Refiz os valores numéricos somando as amplitudes em python3 e eles se confirmam. As dúvidas listadas são de baixa importância e não mudam a nota de nenhum item: o significado de I\_1($\theta$) no ângulo de cada extremo, o pequeno deslocamento dos extremos porque I\_1 varia com $\theta$, e a condição implícita a $\geq$ $\lambda$. \\
10 & 10,0 & A resolução está certa e bate com o gabarito do Poliedro nos dois itens. No item a, as duas dão tg $\theta$ = 2 sen($\pi$/n)/(2 cos($\pi$/n) - $\surd$3), obtida com a mesma lei dos senos no triângulo OAB e o mesmo fechamento $\varphi$ = $\pi$/n. No item b, as duas dão n > 6, ou seja, n $\geq$ 7, sem limite superior. A resolução explica bem por que n = 6 fica de fora: é o caso em que o raio só tangencia o espelho interno. As contas foram conferidas numericamente. Não há divergência; restam só duas observações menores sobre a leitura do enunciado. \\
\bottomrule\end{longtable}%
\section*{Questão 1}%
\href{https://poliedroresolve.sistemapoliedro.com.br/vestibulares/ime/2026/ime-2-fase-dia-2-24-10-2025-2-fase-dia-2/dissertativa-1-fisica-2-fase-dia-2-ime-2026}{Resolução do gabarito}\par\smallskip%
\begin{longtable}{@{}p{0.06\textwidth} p{0.33\textwidth} p{0.33\textwidth} p{0.1\textwidth} p{0.08\textwidth}@{}}\toprule Item & Resolução & Gabarito & Veredito & Certo \\ \midrule
a & {-}2 m < x < 4 m (poço em torno do equilíbrio estável x = 0, com E < 32/3 J; extremos excluídos) & {-}2 < x < 4 (m): a partícula fica presa no poço em torno de x = 0; x = 4 é equilíbrio instável e x = {-}2 tem a mesma Ep que x = 4 & igual & ambos \\
b & E\_sup = Ep(4) = 32/3 J $\approx$ 10,67 J (supremo, não atingido) & E\_sup = 32/3 J ($\approx$ 10,7 J) = Ep(4) = Ep(-2) & igual & ambos \\
c & v(0) = -v0 = -$\surd$(2E/m) (mínimo local); v sobe e se anula no retorno x\_t < -2 em t1 (inflexão, inclinação máxima); chega a +v0 em x = 0 em t2 = 2t1 (máximo local); cai, com inflexão em x = 2, até o mínimo v\_min = $\surd$(2(E-32/3)/m) > 0 em x = 4 (t3); depois cresce sem limite, com assíntota vertical em tempo finito t*. A partícula escapa e não volta. & Começa em -V0, |v| diminui até v = 0 em T1 (retorno em x < 0, inclinação máxima); sobe até +V0 em T2 = 2T1 (x = 0); cai até um mínimo positivo em T3 (x = 4, passa sem parar porque E > E\_sup); depois cresce sem limite, cada vez mais íngreme. & igual & ambos \\
\bottomrule\end{longtable}%
\begin{itemize}%
\item%
\textbf{a)} Mesmo intervalo aberto e mesma justificativa: equilíbrios em F=0 (x=0 estável, x=4 instável) e Ep(x) - 32/3 = -(x-4)$^2$(x+2)/3, então Ep(-2) = Ep(4) = 32/3 J. Conferi a identidade numericamente (erro de 3e{-}14). O enunciado bate com o PDF original.%
\item%
\textbf{b)} Ep(4) = 32 - 64/3 = 32/3 J. A diferença 10,67 contra 10,7 é só arredondamento. Na resolução escrita do gabarito aparece ?2(4)$^2$ - (4(4)$^3$)/3?, com um fator 4 a mais (erro de digitação), mas o resultado final 32/3 J está certo.%
\item%
\textbf{c)} Os dois esboços têm a mesma forma: -v0 $\to$ 0 (T1) $\to$ +v0 (T2 = 2T1) $\to$ mínimo positivo em x = 4 $\to$ crescimento ilimitado. A resolução vai além e detalha concavidades, a inflexão em x = 2 (F?(2) = 0) e o valor v\_min = $\surd$(2(E-32/3)/m). Conferi os números do esboço (m = 1 kg, E = 12 J) com RK4 próprio: x\_t = -2,107 m, t1 = 0,653 s, t2 = 1,307 s com v = 4,899 m/s, t\_i = 1,759 s com v = 3,651 m/s, t3 = 2,613 s com v = 1,633 m/s, v(3,8 s) = 15,78 m/s e t* = t3 + $\int$$_{4}$\^{}$\infty$ dx/v = 4,701 s. Tudo coincide com o texto. A assíntota em tempo finito também está certa: para x grande, v \textasciitilde{} x\^{}\{3/2\}, a integral $\int$dx/v converge e v $\propto$ (t*-t)\^{}\{-3\}. Isso é compatível com o ?cresce sem limite? do gabarito.%
\end{itemize}%
\paragraph{Pontos de dúvida}%
\begin{itemize}%
\item%
\textit{(key)} Na resolução do item b, o gabarito escreve ?E\_sup = 2(4)$^2$ - (4(4)$^3$)/3?, com um fator 4 a mais no termo cúbico. Pela expressão como está escrita, daria 32 - 256/3 < 0. É erro de digitação, porque o resultado declarado (32/3 J) está certo.%
\item%
\textit{(statement)} No item a, ?intervalo de posições possíveis? pode ser lido como intervalo aberto ou fechado. Os dois textos dão o intervalo aberto (-2, 4), que é o certo: com E = 32/3 J a partícula tende assintoticamente a x = 4 e não oscila. Ainda assim, um corretor poderia aceitar {[}-2, 4{]}.%
\item%
\textit{(solution)} No item c, a resolução diz que v $\to$ $\infty$ em tempo finito t* (assíntota vertical). Isso está matematicamente certo para F = x$^2$ - 4x, mas vai além do que o gabarito pede (?cresce sem limite?). Não há conflito: é um detalhe a mais, que a própria resolução marca como idealização.%
\item%
\textit{(solution)} O esboço numérico usa m = 1 kg e E = 12 J, valores que o enunciado não dá (a massa não é fornecida). A resolução deixa claro que o gráfico é qualitativo e vale para qualquer m e qualquer E > 32/3 J.%
\end{itemize}%
\section*{Questão 2}%
\href{https://poliedroresolve.sistemapoliedro.com.br/vestibulares/ime/2026/ime-2-fase-dia-2-24-10-2025-2-fase-dia-2/dissertativa-2-fisica-2-fase-dia-2-ime-2026}{Resolução do gabarito}\par\smallskip%
\begin{longtable}{@{}p{0.06\textwidth} p{0.33\textwidth} p{0.33\textwidth} p{0.1\textwidth} p{0.08\textwidth}@{}}\toprule Item & Resolução & Gabarito & Veredito & Certo \\ \midrule
a & r\_A/r\_B = 16/9 & R\_A/R\_B = 16/9 & igual & ambos \\
b & v\_A/v\_B = 3/4 & v\_A/v\_B = 3/4 & igual & ambos \\
c & M\_B/M\_CA = 1/8 & m\_B/m\_CA = 1/8 & igual & ambos \\
d & v\_B/v\_CA = 2 & v\_B/v\_CA = 2 & igual & ambos \\
e & Apoastro: a = 8/7 r\_B < r\_A = 16/9 r\_B, e = 5/9; além disso v\_CA = 0,5 v\_B < v\_A = 0,75 v\_B (a velocidade é menor que a circular local) & Apoastro: v\_CA < v\_A, então pela vis{-}viva a < R\_A e R\_A = a + c (CA está no ponto mais afastado de S) & igual & ambos \\
\bottomrule\end{longtable}%
\begin{itemize}%
\item%
\textbf{a)} T\_A = 640 d e T\_B = 270 d (dias terrestres). Pela 3ª lei de Kepler, (640/270)$^2$ = (16/9)$^3$ = 4096/729, o que confirma a razão 16/9.%
\item%
\textbf{b)} v = $\surd$(GM\_S/r), então v\_A/v\_B = $\surd$(9/16) = 3/4. A resolução chega ao mesmo valor também por v = 2$\pi$r/T.%
\item%
\textbf{c)} Dividindo as duas razões de forças, (M\_B/M\_CA)(r\_B/r\_A)$^2$ = 81/2048. Então M\_B/M\_CA = (81/2048)$\cdot$(256/81) = 1/8. A resolução e o gabarito fazem a mesma divisão, só que em ordens inversas.%
\item%
\textbf{d)} Com M\_CA = 8M\_B e r\_A = (16/9)r\_B, a condição E\_CA = 7E\_B dá $\frac12$v\_CA$^2$ - (9/16)v\_B$^2$ = -(7/16)v\_B$^2$. Daí v\_CA$^2$ = v\_B$^2$/4 e v\_B/v\_CA = 2. Confirmei com frações exatas.%
\item%
\textbf{e)} A conclusão é a mesma. A resolução vai além do gabarito em dois pontos. Primeiro, justifica por que o ponto do choque é um apside: a colisão é frontal sobre a tangente, então a velocidade fica perpendicular ao raio vetor. O gabarito apenas supõe que 'a colisão ocorre em um extremo'. Segundo, calcula a = 8r\_B/7 e e = 5/9. Conferi que periastro = 32r\_B/63 e que v\_p = 1,75 v\_B dá a mesma energia específica, -(7/16)v\_B$^2$.%
\end{itemize}%
\paragraph{Pontos de dúvida}%
\begin{itemize}%
\item%
\textit{(statement)} Diferença de transcrição. O PDF original diz 'o choque foi perfeitamente inelástico e QUANDO os corpos se moviam em sentido contrário'; a transcrição diz 'e QUE'. A frase é ambígua: pode querer dizer que C e A tinham sentidos opostos antes do choque (choque frontal) ou que CA passou a girar em sentido contrário ao de A. Isso não muda nenhuma resposta, porque só o módulo de v\_CA e a direção tangencial importam.%
\item%
\textit{(key)} No item e, o gabarito supõe sem justificar que o choque acontece num extremo da órbita ('Se a colisão ocorre em um extremo da órbita'). A justificativa completa é a da resolução: num choque frontal, C se move sobre a tangente à órbita de A, então a velocidade de CA é perpendicular ao raio vetor e o ponto é um apside. Numa correção rigorosa, esse argumento pode pesar na nota.%
\item%
\textit{(solution)} A resolução e o gabarito usam hipóteses implícitas razoáveis. A primeira é M\_S ? massa dos planetas, o que permite usar Kepler e v = $\surd$(GM\_S/r) sem massa reduzida. A segunda é tomar a distância de CA a S igual a r\_A logo após o choque, porque os raios dos planetas são desprezíveis. A resolução declara as duas hipóteses.%
\end{itemize}%
\section*{Questão 3}%
\href{https://poliedroresolve.sistemapoliedro.com.br/vestibulares/ime/2026/ime-2-fase-dia-2-24-10-2025-2-fase-dia-2/dissertativa-3-fisica-2-fase-dia-2-ime-2026}{Resolução do gabarito}\par\smallskip%
\begin{longtable}{@{}p{0.06\textwidth} p{0.33\textwidth} p{0.33\textwidth} p{0.1\textwidth} p{0.08\textwidth}@{}}\toprule Item & Resolução & Gabarito & Veredito & Certo \\ \midrule
a & d\_min = 3?/5 & d\_mín = 3?/5 & igual & ambos \\
b & R = ?/14 (o fio afrouxa a 30$^\circ$ acima da horizontal, v$^2$ = 2g?/7, v\_x$^2$ = g?/14) & $\rho$ = ?/14 (o fio afrouxa com cos $\theta$ = 1/2, $\theta$ medido da vertical, v$^2$ = 2g?/7, g$\rho$ = v$^2$ cos$^2$ $\theta$) & igual & ambos \\
\bottomrule\end{longtable}%
\begin{itemize}%
\item%
\textbf{a)} Refiz a conta. Com r = ? - d, o topo do círculo de centro B fica ? - 2d abaixo do ponto de partida. A energia dá v$^2$ = 2g(2d - ?), e no topo com T = 0 vale v$^2$ = g(? - d). Igualando: 5d = 3?, logo d = 3?/5. A resolução chega ao mesmo resultado pela expressão geral T = mg(2d/r - 3 sen $\theta$) aplicada em $\theta$ = 90$^\circ$.%
\item%
\textbf{b)} Refiz a conta. Com r = 4?/7, T = 0 dá sen $\theta$ = 2d/(3r) = 1/2, isto é, 30$^\circ$ acima da horizontal, que é o mesmo que 60$^\circ$ com a vertical no gabarito. Então v$^2$ = g$\cdot$r$\cdot$sen $\theta$ = 2g?/7, e a energia confirma esse valor. No ponto mais alto do lançamento oblíquo a velocidade é v\_x = v/2, logo $\rho$ = v\_x$^2$/g = ?/14. Conferi também numericamente: o ápice fica em (-3$\surd$3/14 ?; 11/28 ?) em relação a B, a 0,540? de B, menos que r = 0,571?. Portanto o fio continua frouxo, a parábola fica dentro do círculo e a esfera não toca a barra antes do ápice.%
\end{itemize}%
\paragraph{Pontos de dúvida}%
\begin{itemize}%
\item%
\textit{(statement)} No item a, o topo da semicircunferência, no ponto (0, r) em relação a B, cai sobre a própria barra, porque r = 2?/5 < d = 3?/5. A esfera chega à barra exatamente ao fim da semicircunferência. Essa é a leitura pretendida, e nem a resolução nem o gabarito comentam o ponto.%
\item%
\textit{(solution)} A resolução e o gabarito medem $\theta$ de referências diferentes: a resolução a partir da horizontal (sen $\theta$ = 1/2), o gabarito a partir da vertical (cos $\theta$ = 1/2). É só convenção, e o ponto de afrouxamento é o mesmo.%
\item%
\textit{(statement)} No item b, 'ponto mais alto' foi entendido como o ápice do lançamento oblíquo depois que o fio afrouxa, e não como o ponto de afrouxamento. A resolução e o gabarito adotam essa leitura, que é a natural.%
\end{itemize}%
\section*{Questão 4}%
\href{https://poliedroresolve.sistemapoliedro.com.br/vestibulares/ime/2026/ime-2-fase-dia-2-24-10-2025-2-fase-dia-2/dissertativa-4-fisica-2-fase-dia-2-ime-2026}{Resolução do gabarito}\par\smallskip%
\begin{longtable}{@{}p{0.06\textwidth} p{0.33\textwidth} p{0.33\textwidth} p{0.1\textwidth} p{0.08\textwidth}@{}}\toprule Item & Resolução & Gabarito & Veredito & Certo \\ \midrule
Q4 --- energia dissipada na malha infinita (0 a 60 s) & E = 7560 J (R\_eq = 2 $\Omega$; P = 50 W, 128 W e 200 W em cada intervalo de 20 s) & E = 7560 J (R\_eq = 2 $\Omega$; 1000 J + 2560 J + 4000 J) & igual & ambos \\
\bottomrule\end{longtable}%
\begin{itemize}%
\item%
\textbf{Q4 --- energia dissipada na malha infinita (0 a 60 s))} Conferi a figura no PDF original. S2 liga o ramo da esquerda ao do meio, S1 liga o do meio ao da direita e S0 liga o ramo da direita à escada, que começa com um R2 em série. Assim, há 1, 2 e 3 ramos em paralelo em cada intervalo, como a resolução diz. R\_eq = R2(1+$\surd$5)/2 = 2 $\Omega$. Os equivalentes de Thévenin são 40 V com 6, 3 e 2 $\Omega$. As correntes são 5, 8 e 10 A, e as potências 50, 128 e 200 W. A energia é 20$\cdot$378 = 7560 J. Refiz as contas em Python e os resultados batem. Também conferi a soma da série na escada, que dá exatamente 1/R\_eq = 0,5 S.%
\end{itemize}%
\paragraph{Pontos de dúvida}%
\begin{itemize}%
\item%
\textit{(statement)} A resolução e o gabarito supõem fontes ideais com a mesma polaridade (a figura mostra os três símbolos de bateria na mesma orientação), chaves ideais e fechamento instantâneo. Supõem também que só se pede a energia na malha infinita, sem contar a dissipada nos R1. Essas hipóteses são razoáveis e estão de acordo com a figura.%
\end{itemize}%
\section*{Questão 5}%
\href{https://poliedroresolve.sistemapoliedro.com.br/vestibulares/ime/2026/ime-2-fase-dia-2-24-10-2025-2-fase-dia-2/dissertativa-5-fisica-2-fase-dia-2-ime-2026}{Resolução do gabarito}\par\smallskip%
\begin{longtable}{@{}p{0.06\textwidth} p{0.33\textwidth} p{0.33\textwidth} p{0.1\textwidth} p{0.08\textwidth}@{}}\toprule Item & Resolução & Gabarito & Veredito & Certo \\ \midrule
a & C = 22/3 $\approx$ 7,3 N; D = 200/3 $\approx$ 66,7 N; E = 26 N; G = 26 N & F\_C = 22/3 N ($\approx$ 7,33 N); F\_D = 200/3 N ($\approx$ 66,7 N); F\_E = 26 N; F\_G = 26 N & igual & ambos \\
b & m\_máx = 8,4 kg (C = 0, G = E = 48 N, D = 96 N) & m = 8,4 kg (iminência de CG tombar em torno de D, F\_C = 0, F\_G = 48 N) & igual & ambos \\
\bottomrule\end{longtable}%
\begin{itemize}%
\item%
\textbf{a)} Refiz com frações exatas. No trecho GE, o peso de GE (12 N) e o do bloco (40 N) atuam no ponto médio x = 9 m, então G = E = 52/2 = 26 N. No trecho CG, tomando momentos em D: 6C + 2$\cdot$26 = 2$\cdot$48, logo C = 22/3 N, e D = 48 + 26 - 22/3 = 200/3 N. As conferências fecham: a soma das reações dá 100 N = (M+m)g, e os momentos em C dão 660 N$\cdot$m dos dois lados. Os dois usam o mesmo modelo: GE apoiado no dente em G, e CG recebendo G para baixo na extremidade em balanço.%
\item%
\textbf{b)} Os dois usam o mesmo critério: a reação em C se anula. Com C = 0, os momentos em D dão 2G = 2$\cdot$48, logo G = 48 N. No trecho GE, 2G = 12 + 10m, o que dá 96 = 12 + 10m e m = 8,4 kg. Pela expressão geral da resolução, C = (42 - 5m)/3 $\geq$ 0, que leva ao mesmo m $\leq$ 8,4 kg. No limite, D = 96 N e E = 48 N (a soma, 144 N, é igual a (6 + 8,4)$\cdot$10). A observação do gabarito de que GE não tomba porque as cargas são simétricas também está correta.%
\end{itemize}%
\paragraph{Pontos de dúvida}%
\begin{itemize}%
\item%
\textit{(statement)} O item (b) fala no maior m 'de forma que a articulação consiga manter a viga em equilíbrio'. A resolução e o gabarito leram isso como a iminência de CG tombar em torno de D (C = 0). Como o enunciado não dá a resistência do dente, essa é a única leitura que dá para calcular, mas a frase poderia sugerir um limite de resistência da própria articulação.%
\item%
\textit{(statement)} O item (a) pede a 'reação de apoio em G', mas G é uma força interna, trocada entre os trechos CG e GE. Os dois respondem com o módulo, 26 N. A resolução diz que essa força é para cima sobre GE e para baixo sobre CG, o que está de acordo com o detalhe do dente na figura e com a observação de que GE 'passa a ser apoiada' em G.%
\item%
\textit{(solution)} A resolução supõe, sem prova completa, que os apoios são unilaterais e que não há forças horizontais. As duas hipóteses são razoáveis para 'apoios verticais simples' e são as mesmas do gabarito. Para o item (a), as frações exatas 22/3 e 200/3 são preferíveis aos valores arredondados 7,3 e 66,7, mas a resolução apresenta as duas formas.%
\end{itemize}%
\section*{Questão 6}%
\href{https://poliedroresolve.sistemapoliedro.com.br/vestibulares/ime/2026/ime-2-fase-dia-2-24-10-2025-2-fase-dia-2/dissertativa-6-fisica-2-fase-dia-2-ime-2026}{Resolução do gabarito}\par\smallskip%
\begin{longtable}{@{}p{0.06\textwidth} p{0.33\textwidth} p{0.33\textwidth} p{0.1\textwidth} p{0.08\textwidth}@{}}\toprule Item & Resolução & Gabarito & Veredito & Certo \\ \midrule
a & i\_min = 1,25 A (em t = 0, v = -2 m/s) & i\_mín = 5/4 A = 1,25 A (v mínima, v = -2 m/s) & igual & ambos \\
b & E = 56,25 J & E = 56,25 J & igual & ambos \\
c & F(5 s) = 2,8125 N (= 45/16 N), para a direita & |F(5)| = 125/16 - 5 = 2,8125 N, para a direita & igual & ambos \\
\bottomrule\end{longtable}%
\begin{itemize}%
\item%
\textbf{a)} Os dois usam $\varepsilon$ = 10 V, r = 2 $\Omega$ e a mesma equação de malha $\varepsilon$ + BLv = (R + r)i, que dá i = 2,5 + 0,625v. No intervalo, v vai de -2 a 2 m/s e i vai de 1,25 a 3,75 A, sempre positiva. O menor módulo é 1,25 A, em t = 0. Conferi as contas em python.%
\item%
\textbf{b)} Com v = 2 m/s, i = (10 + 5)/4 = 3,75 A. Então E = r$\cdot$i$^2$$\cdot$$\Delta$t = 2$\cdot$3,75$^2$$\cdot$2 = 56,25 J. A resolução ainda confere com um balanço de potência: $\varepsilon$i + Fv = 37,5 + 18,75 = 56,25 W, igual a (r + R)i$^2$, e metade vai para r.%
\item%
\textbf{c)} Em t = 5 s, v = 1 m/s e a = -1 m/s$^2$. A corrente é i = 3,125 A e a força magnética vale F\_m = BLi = 7,8125 N, para a esquerda. Pela 2ª lei, F = F\_m + ma = 7,8125 - 5 = 2,8125 N > 0, logo para a direita. 45/16 e 125/16 - 5 são o mesmo número.%
\end{itemize}%
\paragraph{Pontos de dúvida}%
\begin{itemize}%
\item%
\textit{(statement)} Os itens (b) e (c) dependem do sinal da fem de movimento, e esse sinal vem da Figura 1. É preciso saber que A (terminal positivo, já que V\_AB > 0) está ligado ao trilho do fundo, com i para a direita, e que B aponta para cima. Assim a fem BLv se soma a $\varepsilon$ quando v > 0, o que também concorda com Lenz, pois a área do circuito cresce. Conferi no PDF original e a leitura está certa; gabarito e resolução usam o mesmo sinal. Com o sinal oposto, (a) continuaria 1,25 A, mas (b) daria 6,25 J e (c) daria F $\approx$ -0,3125 N, para a esquerda. A resolução registra essa sensibilidade.%
\item%
\textit{(solution)} Algumas hipóteses ficam implícitas, mas são razoáveis e as duas soluções as usam: autoindutância do circuito desprezível (a corrente acompanha v instantaneamente), fonte linear de Thévenin ($\varepsilon$ = 10 V, r = 2 $\Omega$) e corrente dentro da faixa 0--5 A da Figura 2. A resolução as declara explicitamente.%
\end{itemize}%
\section*{Questão 7}%
\href{https://poliedroresolve.sistemapoliedro.com.br/vestibulares/ime/2026/ime-2-fase-dia-2-24-10-2025-2-fase-dia-2/dissertativa-7-fisica-2-fase-dia-2-ime-2026}{Resolução do gabarito}\par\smallskip%
\begin{longtable}{@{}p{0.06\textwidth} p{0.33\textwidth} p{0.33\textwidth} p{0.1\textwidth} p{0.08\textwidth}@{}}\toprule Item & Resolução & Gabarito & Veredito & Certo \\ \midrule
a & V(a) = kQ/(2a) & V$_{1}$(a) = kQ/(2a) & igual & ambos \\
b & Na cavidade, V = kQ/r - kQ/(6a); no metal (2a $\leq$ r $\leq$ 3a), V = kQ/(3a); fora, V = kQ/r. V(a) = 5kQ/(6a), V(2a) = V(3a) = kQ/(3a). Gráfico esboçado. & Para r $\leq$ 2a, V$_{2}$ = kQ(1/r - 1/(6a)); para 2a $\leq$ r $\leq$ 3a, V$_{2}$ = kQ/(3a); para r $\geq$ 3a, V$_{2}$ = kQ/r. V$_{2}$(a) = 5kQ/(6a), V$_{2}$(2a) = V$_{2}$(3a) = kQ/(3a) & igual & ambos \\
c & Casca menor: interna -Q, externa +4Q/5 (carga total -Q/5). Casca maior: interna -Q, externa +6Q/5 (carga total +Q/5). & Casca da Fig. 1: interna -Q, externa +4Q/5. Casca da Fig. 2: interna -Q, externa +6Q/5 & igual & ambos \\
\bottomrule\end{longtable}%
\begin{itemize}%
\item%
\textbf{a)} Os dois usam a mesma distribuição: +Q no centro, -Q em r=a e +Q em r=2a. Por superposição, V(a) = kQ/a - kQ/a + kQ/(2a) = kQ/(2a). A resolução também observa que r=a está na superfície interna do metal, que está no potencial do condutor, kQ/(2a). Conferi com frações exatas.%
\item%
\textbf{b)} As expressões por trecho e os valores destacados são iguais. A resolução tira a constante da continuidade em r=2a: C = kQ/(3a) - kQ/(2a) = -kQ/(6a). O gabarito chega ao mesmo resultado por superposição. Verifiquei: 1 - 1/2 + 1/3 = 5/6. O gráfico da resolução também está certo: hipérbole decrescente na cavidade, patamar no metal e kQ/r fora. Além disso, a resolução espelha a curva para r < 0, o que é coerente com o eixo r da figura, que vai de -3a a 3a.%
\item%
\textbf{c)} Pela lei de Gauss, cada superfície interna fica com -Q. A carga total das cascas se conserva (0), então as externas somam 2Q. Com as cascas distantes, a igualdade de potenciais dá q/(2a) = (2Q - q)/(3a), logo 5q = 4Q e q = 4Q/5. A outra externa fica com 6Q/5. A resolução monta a conta com as cargas totais das cascas (q$_{1}$ = -Q/5, q$_{2}$ = +Q/5), o que é equivalente. Conferi: V = 2kQ/(5a) nas duas cascas. Ao retirar o fio, a distribuição não muda.%
\end{itemize}%
\paragraph{Pontos de dúvida}%
\begin{itemize}%
\item%
\textit{(solution)} No item c, a resolução (e também o gabarito) assume que o fio é muito longo, guarda carga desprezível e não altera os potenciais. Assume ainda que, com as cascas muito distantes, uma não influencia a outra, de modo que cada potencial é k$\cdot$(carga externa)/R\_ext. A hipótese é razoável e vem das observações do enunciado.%
\item%
\textit{(statement)} No item a, r = a fica exatamente na superfície interna da casca da Figura 1. Como o potencial é contínuo ali, o valor kQ/(2a) não muda se o ponto for visto pelo lado da cavidade ou pelo lado do metal.%
\item%
\textit{(solution)} O gráfico do item b foi traçado também para r < 0, como função par, seguindo o eixo da figura de -3a a 3a. Já o gráfico do gabarito, pela transcrição, parece cobrir apenas r $\geq$ 0. Isso não altera nenhum valor pedido.%
\end{itemize}%
\section*{Questão 8}%
\href{https://poliedroresolve.sistemapoliedro.com.br/vestibulares/ime/2026/ime-2-fase-dia-2-24-10-2025-2-fase-dia-2/dissertativa-8-fisica-2-fase-dia-2-ime-2026}{Resolução do gabarito}\par\smallskip%
\begin{longtable}{@{}p{0.06\textwidth} p{0.33\textwidth} p{0.33\textwidth} p{0.1\textwidth} p{0.08\textwidth}@{}}\toprule Item & Resolução & Gabarito & Veredito & Certo \\ \midrule
a & Q\_f = 35000 + 72500 + 12500 = 120000 kcal = 4,8$\cdot$10\^{}8 J; potência frigorífica média $\approx$ 83,3 kW; COP\_Carnot = 246/61,5 = 4, COP = 3; potência do compressor $\approx$ 27,8 kW $\approx$ 37 HP & P $\approx$ 37 HP (Q\_f = 120000 kcal = 480000 kJ; Q\_f = 250/3 kW = 1000/9 HP; $\beta$ = 3; P = Q\_f/$\beta$ $\approx$ 37 HP) & igual & ambos \\
b & 30 HP não bastam (cerca de 19\% abaixo dos \textasciitilde{}37 HP; retiraria só 97200 kcal, ou seja, 81\% do necessário, e o congelamento levaria \textasciitilde{}118,5 min). A estimativa do engenheiro não procede. & A estimativa não procede: são necessários cerca de 37 HP, mais do que os 30 HP indicados. & igual & ambos \\
\bottomrule\end{longtable}%
\begin{itemize}%
\item%
\textbf{a)} Refiz as contas e os valores se confirmam: Q = 1000$\cdot${[}1$\cdot$35 + 72,5 + 0,5$\cdot$25{]} = 120000 kcal = 4,8$\cdot$10\^{}8 J. Em 5760 s isso dá 83,33 kW. $\beta$\_Carnot = 246/(307,5-246) = 4, então $\beta$ = 3. W = 27,78 kW = 37,04 HP. Usar 273,15 K em vez de 273 K dá $\beta$ $\approx$ 3,002 e P $\approx$ 37,01 HP, sem diferença prática. A resolução e o gabarito seguem o mesmo caminho e chegam ao mesmo valor.%
\item%
\textbf{b)} A conclusão é a mesma: 37 HP > 30 HP, logo a estimativa 'com folga' não procede. Conferi as verificações extras da resolução: 3$\cdot$22,5 kW$\cdot$5760 s = 3,888$\cdot$10\^{}8 J = 97200 kcal, que é 81\% de 120000 kcal; o tempo fica 96$\cdot$37,04/30 $\approx$ 118,5 min. As porcentagens também são coerentes entre si: 37/30 $\approx$ 1,23 (23\% acima) e 30/37 $\approx$ 0,81 (19\% abaixo).%
\end{itemize}%
\paragraph{Pontos de dúvida}%
\begin{itemize}%
\item%
\textit{(solution)} A resolução e o gabarito supõem implicitamente que todo o calor retirado pelo evaporador vem do suco. Isso significa câmara adiabática e nenhuma carga de embalagens ou ventiladores. Também supõem potência média constante ao longo dos 96 min e $\beta$ de Carnot calculado com a temperatura da câmara (-27 $^\circ$C) durante todo o processo, embora o suco esteja mais quente no começo. A resolução declara essas hipóteses. É a leitura que o enunciado pede.%
\item%
\textit{(solution)} A resolução diz que 118,5 min é 'cerca de 2 h'. Na verdade são aproximadamente 1 h 58 min, então é só um arredondamento de redação, e esse cálculo extra não é exigido pelo item b.%
\end{itemize}%
\section*{Questão 9}%
\href{https://poliedroresolve.sistemapoliedro.com.br/vestibulares/ime/2026/ime-2-fase-dia-2-24-10-2025-2-fase-dia-2/dissertativa-9-fisica-2-fase-dia-2-ime-2026}{Resolução do gabarito}\par\smallskip%
\begin{longtable}{@{}p{0.06\textwidth} p{0.33\textwidth} p{0.33\textwidth} p{0.1\textwidth} p{0.08\textwidth}@{}}\toprule Item & Resolução & Gabarito & Veredito & Certo \\ \midrule
a & I($\theta$) = I\_1($\theta$){[}4 + 2$\surd$3 cos(2$\pi$ a sen$\theta$/$\lambda$){]} & I\_R($\theta$) = I\_1($\theta$){[}4 + 2$\surd$3 cos((2$\pi$/$\lambda$) a sen$\theta$){]} = 2I\_1($\theta$){[}2 + $\surd$3 cos((2$\pi$a/$\lambda$) sen$\theta$){]} & igual & ambos \\
b & I\_máx = (4+2$\surd$3) I\_1($\theta$) = (1+$\surd$3)$^2$ I\_1 $\approx$ 7,46 I\_1, em sen$\theta$ = $\lambda$/a & I\_1$^\circ$máx = I\_1($\theta$)(4 + 2$\surd$3), em sen$\theta$ = $\lambda$/a & igual & ambos \\
c & I\_mín = (4-2$\surd$3) I\_1($\theta$) = ($\surd$3-1)$^2$ I\_1 $\approx$ 0,54 I\_1, em sen$\theta$ = $\lambda$/(2a) & I\_1$^\circ$mín = I\_1($\theta$)(4 - 2$\surd$3), em sen$\theta$ = $\lambda$/(2a) & igual & ambos \\
\bottomrule\end{longtable}%
\begin{itemize}%
\item%
\textbf{a)} A expressão é a mesma. Ambos usam I\_2 = 3I\_1, a fórmula I = I\_1 + I\_2 + 2$\surd$(I\_1I\_2)cos$\Delta$? e $\Delta$? = 2$\pi$a sen$\theta$/$\lambda$. Pela figura, F\_1 está acima e P também, então r\_2 - r\_1 = a sen$\theta$ > 0. Como só entra cos$\Delta$?, o sinal da diferença de fase não altera o resultado.%
\item%
\textbf{b)} O valor e a posição são os mesmos. Conferi numericamente: |A\_1 + A\_2 e\^{}\{i$\cdot$2$\pi$\}|$^2$, com A\_2 = $\surd$3$\cdot$A\_1, dá (1+$\surd$3)$^2$ = 4+2$\surd$3 $\approx$ 7,464, o que confirma o resultado.%
\item%
\textbf{c)} O valor e a posição são os mesmos. Conferi numericamente: ($\surd$3-1)$^2$ = 4-2$\surd$3 $\approx$ 0,536. O mínimo não é nulo porque as amplitudes das duas fontes são diferentes.%
\end{itemize}%
\paragraph{Pontos de dúvida}%
\begin{itemize}%
\item%
\textit{(statement)} I\_1($\theta$) depende de $\theta$. Nos itens b e c, o resultado 'em função de I\_1($\theta$)' vale com I\_1 calculado no ângulo do extremo, isto é, $\theta$ = arcsen($\lambda$/a) ou $\theta$ = arcsen($\lambda$/2a). A resolução diz isso explicitamente no item b. O gabarito deixa só implícito.%
\item%
\textit{(solution)} Resolução e gabarito localizam os extremos por cos$\Delta$? = $\pm$1, como se I\_1 fosse constante. Como I\_1($\theta$) varia com $\theta$, a posição exata dos extremos locais de I($\theta$) se desloca um pouco. A resolução declara essa hipótese de forma explícita. É a leitura esperada pela banca.%
\item%
\textit{(statement)} O primeiro máximo com $\theta$ > 0 só existe se a $\geq$ $\lambda$, porque sen$\theta$ = $\lambda$/a $\leq$ 1. Da mesma forma, o primeiro mínimo exige a $\geq$ $\lambda$/2. Nem a resolução nem o gabarito mencionam essa condição. Ela é implícita no contexto de Young.%
\end{itemize}%
\section*{Questão 10}%
\href{https://poliedroresolve.sistemapoliedro.com.br/vestibulares/ime/2026/ime-2-fase-dia-2-24-10-2025-2-fase-dia-2/dissertativa-10-fisica-2-fase-dia-2-ime-2026}{Resolução do gabarito}\par\smallskip%
\begin{longtable}{@{}p{0.06\textwidth} p{0.33\textwidth} p{0.33\textwidth} p{0.1\textwidth} p{0.08\textwidth}@{}}\toprule Item & Resolução & Gabarito & Veredito & Certo \\ \midrule
a & tan $\theta$ = sen($\pi$/n)/(cos($\pi$/n) - $\surd$3/2) = 2 sen($\pi$/n)/(2 cos($\pi$/n) - $\surd$3) & tg $\theta$ = 2$\cdot$sen($\pi$/n)/(2$\cdot$cos($\pi$/n) - $\surd$3) & igual & ambos \\
b & n inteiro com n $\geq$ 7 (n > 6), sem limite superior; n = 6 corresponde só à tangência ($\theta$ = 90$^\circ$) & n > 6 (n natural, ou seja, n $\geq$ 7) & igual & ambos \\
\bottomrule\end{longtable}%
\begin{itemize}%
\item%
\textbf{a)} A resolução escreve a mesma expressão do gabarito, inclusive na mesma forma. O caminho também é o mesmo: no triângulo OAB o ângulo em B é $\pi$-$\theta$, a lei dos senos dá sen $\alpha$ = (r/R) sen $\theta$, o ângulo central é $\varphi$ = $\theta$-$\alpha$, e o fechamento com 2n segmentos congruentes dá 2n$\varphi$ = 2$\pi$, ou seja, $\varphi$ = $\pi$/n. Conferi numericamente para n = 7, 8, 12 e 100 que sen($\theta$-$\pi$/n) - ($\surd$3/2) sen $\theta$ = 0 até a precisão da máquina. Também conferi a verificação da própria resolução para n = 12: $\theta$ $\approx$ 68,9$^\circ$ e $\alpha$ $\approx$ 53,9$^\circ$.%
\item%
\textbf{b)} Mesma conclusão e mesmo argumento: tg $\theta$ > 0 exige cos($\pi$/n) > $\surd$3/2, logo $\pi$/n < $\pi$/6 e n > 6. A resolução acrescenta uma leitura física correta. O maior ângulo central possível de um segmento ocorre na tangência e vale arccos(r/R) = $\pi$/6. Com n = 6 o raio só tangencia o espelho interno ($\theta$ $\to$ 90$^\circ$; numericamente, $\theta$ $\approx$ 89,9995$^\circ$ para n = 6,0001), e com n $\leq$ 5 o raio nem chega ao cilindro interno. Para n $\geq$ 7 vale $\alpha$ = $\theta$ - $\pi$/n > 0, o que deixa o triângulo consistente.%
\end{itemize}%
\paragraph{Pontos de dúvida}%
\begin{itemize}%
\item%
\textit{(statement)} O caso degenerado n = 1 (sen $\pi$ = 0, logo tg $\theta$ = 0, ou seja, raio radial indo e voltando entre A e B) não é tratado explicitamente pelo gabarito. A resolução o exclui de passagem com 'sen($\pi$/n) > 0 para n $\geq$ 2'. Fisicamente esse raio não dá uma volta, então excluí{-}lo é razoável. Mesmo assim, como o enunciado diz só n $\in$ $\mathbb{N}$, vale mencionar.%
\item%
\textit{(statement)} A frase 'os raios se sobrepõem após uma volta' foi lida como fechamento após exatamente uma volta (2n$\varphi$ = 2$\pi$). Fechamentos após k voltas (trajetórias em estrela, 2n$\varphi$ = 2$\pi$k) ficam excluídos por essa leitura. Tanto o gabarito quanto a resolução adotam essa interpretação, que é a natural.%
\end{itemize}%
\end{document}